\documentclass[12pt,a4paper]{article}
\usepackage[T2A]{fontenc}
\usepackage[utf8]{inputenc}
\usepackage[russian]{babel}
\usepackage{amsmath,amssymb}
\usepackage{geometry}
\geometry{margin=2.2cm}

\begin{document}

\section*{Грамматика продвинутого лямбда-интерпретатора}

\subsection*{Требования к языку}

Входная программа состоит из набора именованных определений и одного
лямбда-выражения. Каждое определение записывается на отдельной
непустой строке и имеет вид
\begin{verbatim}
let name = expression
\end{verbatim}
и могут использоваться в последующем выражении как обычные переменные.
Слово \verb|let| является ключевым словом и не может быть именем
переменной или определения.

Лямбда обозначается символом \verb|\|. Должны поддерживаться
абстракции с несколькими параметрами:
\begin{verbatim}
\x y z. expression
\end{verbatim}
Такая запись является сокращением для вложенных абстракций:
\[
\lambda x.\lambda y.\lambda z.\ expression.
\]

\subsection*{Лексические правила}

Пробельные символы между токенами игнорируются. Идентификатор начинается
с буквы, далее могут идти буквы, цифры, апостроф и символ подчеркивания.
Ключевое слово \verb|let| исключается из множества идентификаторов.

\begin{verbatim}
letter      ::= "a" | ... | "z" | "A" | ... | "Z"
digit       ::= "0" | ... | "9"
identifier  ::= letter (letter | digit | "'" | "_")*
\end{verbatim}

\subsection*{Факторизованная грамматика}

Грамматика записана в EBNF. Оператор $+$ означает <<один или больше>>,
оператор $*$ используется только в лексическом правиле идентификатора.
В синтаксической части нет явных $\varepsilon$-продукций.

\begin{verbatim}
program           ::= definition* expression
definition        ::= "let" identifier "=" expression lineEnd
expression        ::= abstraction | applicationChain
abstraction       ::= "\" identifier+ "." expression
applicationChain  ::= atom+
atom              ::= identifier | "(" expression ")"
lineEnd           ::= end of current non-empty line
\end{verbatim}

\subsection*{Почему грамматика подходит для парсера}

В исходной грамматике аппликация естественно записывается как
\[
term ::= term\ term,
\]
что дает левую рекурсию. Для FParsec и рекурсивного спуска такая форма
неподходящая: парсер \verb|term| немедленно вызывает сам себя без
потребления входа.

В предложенной грамматике аппликация заменена на цепочку атомов:
\[
applicationChain ::= atom^+.
\]
Такой нетерминал парсится как \verb|many1 atom|. Если в цепочке один
атом, результатом является сам атом. Если атомов несколько, AST строится
левоассоциативно:
\[
a\ b\ c \equiv ((a\ b)\ c).
\]

Альтернативы в правиле \verb|expression| различаются первым токеном:
абстракция начинается с \verb|\|, а цепочка аппликации начинается с
идентификатора или открывающей скобки. Поэтому правило удобно реализовать
через \verb|choice| без попыток угадывать одинаковые префиксы.

Разделение определений по строкам нужно для однозначности. Без этого
запись
\begin{verbatim}
let K = \x y.x
S K K
\end{verbatim}
можно было бы прочитать как одно определение с телом \verb|\x y.x S K K|.
Поэтому на уровне программы непустые строки, начинающиеся с ключевого
слова \verb|let|, считаются определениями, а оставшиеся строки образуют
итоговое выражение.

\subsection*{Построение AST}

Для дальнейшей реализации удобно использовать тот же базовый AST, что и в
предыдущей задаче:
\[
\begin{array}{rcl}
Expr & ::= & Var(name) \\
     & \mid & Abs(name,\ Expr) \\
     & \mid & App(Expr,\ Expr)
\end{array}
\]

Именованные определения могут храниться отдельно в таблице
\[
name \mapsto Expr.
\]
Перед нормализацией входного выражения имена из таблицы можно раскрыть
или учитывать их при редукции как заранее заданные переменные.

Абстракция с несколькими параметрами строится как последовательность
вложенных \verb|Abs|. Например:
\begin{verbatim}
\x y.x  ->  Abs(x, Abs(y, Var(x)))
\end{verbatim}

Цепочка аппликации строится левым сворачиванием. Например:
\begin{verbatim}
S K K  ->  App(App(Var(S), Var(K)), Var(K))
\end{verbatim}

\subsection*{Пример 1}

Вход:
\begin{verbatim}
let S = \x y z.x z (y z)
let K = \x y.x

S K K
\end{verbatim}

Разбор определений:
\[
\begin{array}{rcl}
S & \mapsto & Abs(x,\ Abs(y,\ Abs(z,\ App(App(Var(x), Var(z)),
                 App(Var(y), Var(z)))))) \\
K & \mapsto & Abs(x,\ Abs(y,\ Var(x)))
\end{array}
\]

Разбор итогового выражения:
\begin{verbatim}
S K K  ->  App(App(Var(S), Var(K)), Var(K))
\end{verbatim}

После раскрытия определений и бета-редукции ожидаемый результат:
\begin{verbatim}
\x.x
\end{verbatim}
или любой альфа-эквивалентный терм.

\subsection*{Пример 2}

Вход:
\begin{verbatim}
let I = \x.x

(\x.x) I
\end{verbatim}

Разбор итогового выражения:
\begin{verbatim}
(\x.x) I  ->  App(Abs(x, Var(x)), Var(I))
\end{verbatim}

После раскрытия определения \verb|I| и бета-редукции:
\begin{verbatim}
\x.x
\end{verbatim}

\end{document}
